\section{Related Work}
\label{sec:related}

\paragraph{Computational Argumentation.}
Project Debater~\cite{slonim2021debater} demonstrated AI debate in a
simplified format designed for lay audiences, but explicitly acknowledged
its inability to engage in expert competitive debate. Prior work in
computational argumentation has focused on argument mining, stance detection,
and persuasion prediction rather than end-to-end debate generation, and
large-scale resources such as OpenDebateEvidence~\cite{roush2024opendebateevidence}
have primarily supported evidence retrieval and summarization rather than
trained debaters. We target the International Public Debate Association
(IPDA) format as an instance of competitive debate more broadly. Competitive
debate formats share typed argument structures
(Advantage = Uniqueness + Link + Impact~\cite{snider2008code}), explicit
rules, and specialist judges; we argue that these structured competitive
formats provide the discrete rewards needed for optimization.

\paragraph{Self-Play and Search in Games.}
Our branching mechanism draws on the tradition of self-play and tree search
in game-playing AI. AlphaGo Zero~\cite{silver2017mastering} and
AlphaZero~\cite{silver2018general} achieve superhuman performance via
Monte Carlo tree search (MCTS) guided by learned value and policy networks,
while OpenAI Five~\cite{berner2019dota} scaled self-play to the imperfect-information
setting of Dota~2. Unlike MCTS, which exhaustively explores game trees,
our branching is triggered only when candidate speeches would elicit
different opponent responses, focusing computation on informationally
rich divergence points. This threshold-based branching is closer to
selective search in that it expands only the nodes where counterfactual
outcomes are most likely to differ.

\paragraph{Tree Search with LLMs.}
Similar branching ideas appear in other RL-for-generation domains. Tree of
Thoughts~\cite{yao2023tree} expands reasoning trees at inference time using
self-evaluation heuristics, and RAP~\cite{hao2023reasoning} treats an LLM as
a world model and uses MCTS to search over reasoning steps guided by value
estimates. Our branching is related but operates on a different signal: we
branch based on \emph{opponent-response divergence} during adversarial
multi-agent rollouts rather than on self-consistency or value estimates, and
the resulting winner-shift events become training targets for offline GRPO
rather than an inference-time search procedure.

\paragraph{RL with Verifiable Rewards.}
Recent work on reasoning-oriented RL trains with programmatically verifiable
reward signals. DeepSeek-R1~\cite{guo2025deepseekr1} demonstrates that
large-scale RL with rule-based rewards elicits strong chain-of-thought
behaviors without supervised warm-start. Our three-layer reward hierarchy
occupies a middle ground: speech-level rubrics and branch-based winner-shift
detection provide near-verifiable local signals, while retrospective
multi-model panels handle the non-verifiable strategic quality of full
debates.

\paragraph{Debate as AI Alignment.}
Irving et al.~\cite{irving2018debate} proposed debate as a scalable oversight
mechanism, where competing agents help judges identify truthful answers.
Michael et al.~\cite{michael2023debate} showed that debate achieves 84\%
judge accuracy versus 74\% for consultancy, with debate errors decreasing
with debater skill. Du et al.~\cite{du2024multiagent} demonstrated that
multi-agent debate improves factuality through convergence dynamics. Our
work complements this line by studying the \emph{training dynamics} of
debate models rather than debate as an oversight protocol.

\paragraph{Constitutional AI and RLAIF.}
RLAIF~\cite{lee2024rlaif} replaces human annotators with LLM judges at
25--125$\times$ lower cost. Constitutional AI~\cite{bai2022constitutional}
showed that models can self-critique and revise outputs against a set of
principles, enabling alignment without human preference labels. Our
hierarchical reward system extends this idea: rather than a single
constitution, we employ speech-specific rubrics at Layer~1, causal outcome
attribution at Layer~2, and multi-model retrospective judging at Layer~3,
each providing a different form of AI-generated feedback.

\paragraph{Preference Optimization.}
We build on ORPO~\cite{hong2024orpo}, which eliminates
the reference model requirement of DPO~\cite{rafailov2023dpo}, and
offline GRPO~\cite{shao2024deepseekmath}, which uses importance sampling
ratios $\pi_\theta / \pi_{\text{ref}}$ to normalize for response
length---critical for tasks where quality does not correlate with verbosity.
Our contribution is the \emph{group-wise} application: decomposing
heterogeneous calls into four modality groups and training sequentially
with SFT-on-golden interleaving prevents the cross-modality interference
we observe when training all call types jointly.

\paragraph{LLM-as-Judge.}
LLM judges achieve 80\%+ human agreement but exhibit systematic
biases~\cite{zheng2023mtbench}: position bias, verbosity preference, and
self-enhancement. We document a novel bias pattern specific to adversarial
evaluation: different judge models exhibit \emph{opposing side biases}
(Section~\ref{sec:judge_bias}), motivating multi-model panels for debate
evaluation.
